\section{Mathematical Model and Theory}
\label{sec:mathematical-model}

The reference first formulates the food chain in dimensional variables and
then rescales concentration and time.  Nutrient concentration is measured
relative to the input concentration, time is measured in units of the inverse
chemostat dilution rate, and the organism concentrations are scaled using the
corresponding yield constants.  This nondimensionalization fixes the nutrient
input and dilution terms at one and reduces the model to four state variables
and three effective removal rates.  The model used in the analysis and in our
implementation is
\begin{equation}
\label{eq:model}
\begin{aligned}
\dot S &= 1-S-f_1(S)x,\\
\dot x &= x\bigl(f_1(S)-D_1\bigr)-f_2(x)y,\\
\dot y &= y\bigl(f_2(x)-D_2\bigr)-f_3(y)z,\\
\dot z &= z\bigl(f_3(y)-D_3\bigr).
\end{aligned}
\end{equation}
Here and below a dot denotes differentiation with respect to dimensionless
time.

\subsection{State Variables and Model Parameters}

The variables form the state vector
\(
X(t)=(S(t),x(t),y(t),z(t))^{\mathsf T}\in\mathbb{R}_+^4
\).
Their meanings are summarized in Table~\ref{tab:model-variables}.

\begin{table}[htbp]
\centering
\caption{State variables of the nondimensional chemostat model.}
\label{tab:model-variables}
\begin{tabular}{cll}
\toprule
Symbol & Biological meaning & Domain \\
\midrule
\(S(t)\) & Nutrient concentration & \(S\geq 0\) \\
\(x(t)\) & Prey concentration; consumes the nutrient & \(x\geq 0\) \\
\(y(t)\) & First-predator concentration; consumes \(x\) & \(y\geq 0\) \\
\(z(t)\) & Top-predator concentration; consumes \(y\) & \(z\geq 0\) \\
\bottomrule
\end{tabular}
\end{table}

The paper allows general response functions \(f_i\).  It assumes that every
\(f_i:\mathbb{R}_+\rightarrow\mathbb{R}_+\) is continuously differentiable,
\(f_i(0)=0\), and \(f_i'(u)>0\) for \(u>0\).  Thus, growth increases with the
available food concentration.  For reproducible simulations, the project
specializes these functions to the Monod form
\begin{equation}
\label{eq:monod}
f_i(u)=\frac{a_i u}{k_i+u},\qquad a_i>0,\quad k_i>0,qquad i=1,2,3.
\end{equation}
The parameter \(a_i\) is the maximum specific growth rate and \(k_i\) is the
half-saturation constant, since \(f_i(k_i)=a_i/2\).  The derivative and inverse
needed in the equilibrium calculations are
\begin{equation}
\label{eq:monod-derivative-inverse}
f_i'(u)=\frac{a_i k_i}{(k_i+u)^2},
\qquad
f_i^{-1}(d)=\frac{d k_i}{a_i-d}\quad (0\leq d<a_i).
\end{equation}

Each \(D_i>0\) is the dimensionless total removal rate of one organism.  In
the dimensional model this rate is the sum of chemostat washout and the
organism's specific death rate; division by the dilution rate produces the
dimensionless \(D_i\).  The nine numerical parameters used by the simulator
are therefore
\[
(a_1,k_1,a_2,k_2,a_3,k_3,D_1,D_2,D_3).
\]

\subsection{Governing Differential Equations}

The first equation of~\eqref{eq:model} balances normalized nutrient inflow,
nutrient washout, and consumption by the prey.  The prey grows at per-capita
rate \(f_1(S)\), is removed at rate \(D_1\), and is consumed by the first
predator through the flux \(f_2(x)y\).  Similarly, \(y\) grows through prey
consumption, is removed at rate \(D_2\), and is consumed by \(z\).  The top
predator grows at rate \(f_3(y)\) and is removed at rate \(D_3\).

This chain structure is important: \(y\) does not consume nutrient directly,
and \(z\) consumes neither nutrient nor prey.  The transfer terms cancel when
all four equations are added.  Distinct removal rates prevent the stronger
conservation law available in equal-removal chemostat models, so the complete
four-dimensional system must be studied.

\subsection{Initial Conditions and Parameter Restrictions}

The reference considers strictly positive initial data,
\begin{equation}
\label{eq:initial-condition}
S(0)=S_0>0,\qquad x(0)=x_0>0,\qquad
y(0)=y_0>0,\qquad z(0)=z_0>0.
\end{equation}
Boundary initial states are also meaningful when studying washout equilibria.
All model parameters are positive.  In addition, a finite positive solution
of \(f_i(u)=D_i\) requires \(D_i<a_i\) for the Monod functions.  Stronger
conditions are needed for an equilibrium to lie in the biologically feasible
region; these are derived below rather than imposed globally.

\subsection{Nonnegativity, Boundedness, and Biological Feasibility}

The nonnegative orthant is positively invariant.  At \(S=0\), the first
equation gives \(\dot S=1>0\).  At \(x=0\), \(y=0\), or \(z=0\), the
corresponding derivative is zero because \(f_i(0)=0\).  Therefore a solution
starting with nonnegative concentrations cannot cross into negative values.

To establish boundedness, define the total dimensionless concentration
\[
T(t)=S(t)+x(t)+y(t)+z(t).
\]
Adding the equations in~\eqref{eq:model} gives
\begin{equation}
\label{eq:total-biomass}
\dot T=1-S-D_1x-D_2y-D_3z.
\end{equation}
Let
\[
D_{\min}=\min\{1,D_1,D_2,D_3\},\qquad
D_{\max}=\max\{1,D_1,D_2,D_3\}.
\]
Since all state variables are nonnegative,
\begin{equation}
1-D_{\max}T\leq \dot T\leq 1-D_{\min}T.
\end{equation}
The scalar comparison theorem consequently yields
\begin{equation}
\frac{1}{D_{\max}}
\leq \liminf_{t\to\infty}T(t)
\leq \limsup_{t\to\infty}T(t)
\leq \frac{1}{D_{\min}}.
\end{equation}
Equivalently, for any sufficiently small \(\kappa>0\), trajectories are
eventually attracted to the bounded region
\begin{equation}
\label{eq:absorbing-region}
\Gamma_\kappa=
\left\{X\in\mathbb{R}_+^4:
\frac{1}{D_{\max}}-\kappa\leq S+x+y+z
\leq\frac{1}{D_{\min}}+\kappa\right\}.
\end{equation}
Thus the model is dissipative: long-term trajectories enter a common bounded
set.  This result is both biologically necessary and numerically useful,
because an unbounded simulated trajectory would indicate an invalid parameter
or integration failure rather than intended model behavior.

\subsection{Equilibrium Points}

An equilibrium \(P=(S^*,x^*,y^*,z^*)\) satisfies
\begin{equation}
\label{eq:equilibrium-system}
\begin{aligned}
1-S^*-f_1(S^*)x^*&=0,\\
x^*\bigl(f_1(S^*)-D_1\bigr)-f_2(x^*)y^*&=0,\\
y^*\bigl(f_2(x^*)-D_2\bigr)-f_3(y^*)z^*&=0,\\
z^*\bigl(f_3(y^*)-D_3\bigr)&=0.
\end{aligned}
\end{equation}
The four possible trophic regimes follow by deciding which organism
concentrations are zero.

\subsubsection{Washout Equilibrium \(P_0\)}

If all organisms are absent, nutrient converges to its normalized input level:
\begin{equation}
P_0=(1,0,0,0).
\end{equation}
This equilibrium always exists.

\subsubsection{Prey-Only Equilibrium \(P_1\)}

For \(x_1>0\) and \(y_1=z_1=0\), the prey equation requires
\(f_1(S_1)=D_1\).  The nutrient balance then gives
\begin{equation}
\label{eq:p1}
P_1=(S_1,x_1,0,0),\qquad
S_1=f_1^{-1}(D_1),\qquad
x_1=\frac{1-S_1}{D_1}.
\end{equation}
For Monod growth,
\begin{equation}
S_1=\frac{D_1k_1}{a_1-D_1}.
\end{equation}
The biological condition \(0<S_1<1\) is equivalent to
\(D_1<f_1(1)=a_1/(k_1+1)\).  This is precisely the condition that prey can
increase when rare at \(P_0\).

\subsubsection{Two-Species Equilibrium \(P_2\)}

For prey and the first predator to persist while \(z=0\), the \(y\)-equation
requires \(f_2(x_2)=D_2\).  Therefore
\begin{equation}
\label{eq:p2}
\begin{aligned}
P_2&=(S_2,x_2,y_2,0),\\
x_2&=f_2^{-1}(D_2),\\
S_2+f_1(S_2)x_2&=1,\\
y_2&=\frac{x_2\bigl(f_1(S_2)-D_1\bigr)}{D_2}.
\end{aligned}
\end{equation}
For the Monod function, \(x_2=D_2k_2/(a_2-D_2)\).  Once \(x_2\) is known,
the equation for \(S_2\) has a unique root in \((0,1)\), because
\(S+f_1(S)x_2\) is strictly increasing from zero to a value greater than one.
Finally, \(y_2>0\) requires \(f_1(S_2)>D_1\).  In the equilibrium hierarchy,
this condition is equivalent to the first-predator invasion threshold
\(f_2(x_1)>D_2\).

\subsubsection{Coexistence Equilibrium \(P_3\)}

At an interior equilibrium all four components are positive.  The top-predator
equation first fixes
\begin{equation}
\label{eq:y3}
y_3=f_3^{-1}(D_3)
=\frac{D_3k_3}{a_3-D_3}
\end{equation}
for Monod growth.  The remaining coordinates satisfy
\begin{equation}
\label{eq:p3}
\begin{aligned}
P_3&=(S_3,x_3,y_3,z_3),\\
S_3+f_1(S_3)x_3&=1,\\
x_3\bigl(f_1(S_3)-D_1\bigr)&=f_2(x_3)y_3,\\
z_3&=\frac{y_3\bigl(f_2(x_3)-D_2\bigr)}{D_3}.
\end{aligned}
\end{equation}
The coupled equations for \(S_3\) and \(x_3\) generally do not have a useful
closed form, even after choosing Monod functions.  The implementation solves
\(S+f_1(S)x=1\) by bisection for each trial \(x\), then locates a sign change
in the prey-balance residual.  The coexistence state is feasible only when all
coordinates are positive.  In particular, \(D_3<a_3\) is needed to define
\(y_3\), and \(f_2(x_3)>D_2\) is needed for \(z_3>0\).

\subsection{Existence and Positivity Conditions}

The equilibria form a trophic hierarchy controlled by invasion rates.  An
invasion rate is the per-capita growth rate of a missing species when that
species is introduced at very low concentration.  The three relevant rates
are
\begin{equation}
\label{eq:invasion-rates}
\lambda_x=f_1(1)-D_1,\qquad
\lambda_y=f_2(x_1)-D_2,\qquad
\lambda_z=f_3(y_2)-D_3.
\end{equation}

\begin{table}[htbp]
\centering
\caption{Equilibrium hierarchy and biological existence conditions.}
\label{tab:equilibrium-existence}
\begin{tabular}{lll}
\toprule
Equilibrium & Persisting organisms & Main feasibility condition \\
\midrule
\(P_0\) & None & Always exists \\
\(P_1\) & \(x\) & \(\lambda_x>0\) \\
\(P_2\) & \(x,y\) & \(P_1\) exists and \(\lambda_y>0\) \\
\(P_3\) & \(x,y,z\) & \(P_2\) exists and \(\lambda_z>0\) \\
\bottomrule
\end{tabular}
\end{table}

Positive invasion means that the next trophic level can enter the lower-level
equilibrium.  Negative invasion means that it washes out when rare.  At a zero
invasion rate the corresponding equilibrium is nonhyperbolic, and a change in
the qualitative survival regime may occur.

\subsection{Local Stability and Threshold Conditions}

Linearizing~\eqref{eq:model} at a state \((S,x,y,z)\) gives the Jacobian
\begin{equation}
\label{eq:jacobian}
J=\begin{pmatrix}
-1-f_1'(S)x & -f_1(S) & 0 & 0\\
f_1'(S)x & f_1(S)-D_1-f_2'(x)y & -f_2(x) & 0\\
0 & f_2'(x)y & f_2(x)-D_2-f_3'(y)z & -f_3(y)\\
0 & 0 & f_3'(y)z & f_3(y)-D_3
\end{pmatrix}.
\end{equation}
An equilibrium is locally asymptotically stable when every eigenvalue of its
Jacobian has negative real part.

\paragraph{Stability of \(P_0\).}
At washout, the eigenvalues are
\begin{equation}
-1,\qquad f_1(1)-D_1,\qquad -D_2,\qquad -D_3.
\end{equation}
Consequently,
\begin{equation}
\label{eq:p0-stability}
P_0\text{ is locally asymptotically stable}
\quad\Longleftrightarrow\quad f_1(1)-D_1<0.
\end{equation}
For Monod growth this becomes
\(a_1/(k_1+1)<D_1\).  If the inequality is reversed, prey can invade and
\(P_0\) is unstable in the \(x\)-direction.

\paragraph{Stability of \(P_1\).}
The \((S,x)\) block of \(J(P_1)\) is
\[
J_{Sx}(P_1)=
\begin{pmatrix}
-1-x_1f_1'(S_1)&-D_1\\
x_1f_1'(S_1)&0
\end{pmatrix}.
\]
Its trace is negative and its determinant is positive, so both eigenvalues
have negative real part.  The remaining eigenvalues are
\(
f_2(x_1)-D_2
\)
and \(-D_3\).  Hence
\begin{equation}
\label{eq:p1-stability}
P_1\text{ is locally asymptotically stable}
\quad\Longleftrightarrow\quad f_2(x_1)-D_2<0.
\end{equation}
When this threshold is positive, \(P_1\) is a saddle that repels trajectories
in the \(y\)-direction.

\paragraph{Stability of \(P_2\).}
For a compact Routh--Hurwitz representation, define
\begin{align*}
a&=1+x_2f_1'(S_2), & b&=f_1(S_2), & c&=x_2f_1'(S_2),\\
d&=y_2f_2'(x_2), & e&=D_2, &
r&=f_1(S_2)-D_1-y_2f_2'(x_2).
\end{align*}
The active \((S,x,y)\) block is
\[
\begin{pmatrix}
-a&-b&0\\
c&r&-e\\
0&d&0
\end{pmatrix},
\]
whose characteristic polynomial is
\begin{equation}
\lambda^3+A_1\lambda^2+A_2\lambda+A_3,
\end{equation}
with
\begin{equation}
\label{eq:p2-cubic-coefficients}
A_1=a-r,\qquad
A_2=ed-ar+bc,\qquad
A_3=aed.
\end{equation}
The fourth eigenvalue is the top-predator invasion rate
\(\lambda_z=f_3(y_2)-D_3\).  The cubic Routh--Hurwitz criterion gives
\begin{equation}
\label{eq:p2-stability}
A_1>0,\quad A_2>0,\quad A_3>0,\quad A_1A_2>A_3,
\quad\text{and}\quad f_3(y_2)-D_3<0
\end{equation}
as sufficient and necessary local stability conditions.  If the active block
is stable but \(f_3(y_2)-D_3>0\), then \(P_2\) is a saddle repelling in the
\(z\)-direction.

\paragraph{Stability of \(P_3\).}
At coexistence, introduce the abbreviations
\begin{align*}
A&=1+x_3f_1'(S_3), & B&=f_1(S_3), & C&=x_3f_1'(S_3),\\
R&=f_1(S_3)-D_1-y_3f_2'(x_3),
& E&=f_2(x_3), & F&=y_3f_2'(x_3),\\
Q&=f_2(x_3)-D_2-z_3f_3'(y_3),
& G&=f_3(y_3), & H&=z_3f_3'(y_3).
\end{align*}
Then
\begin{equation}
J(P_3)=
\begin{pmatrix}
-A&-B&0&0\\
C&R&-E&0\\
0&F&Q&-G\\
0&0&H&0
\end{pmatrix}.
\end{equation}
Set \(p=A-R\) and \(q=BC-AR\).  Expanding
\(\det(\lambda I-J(P_3))\) gives
\begin{equation}
\lambda^4+C_1\lambda^3+C_2\lambda^2+C_3\lambda+C_4,
\end{equation}
where
\begin{equation}
\label{eq:p3-quartic-coefficients}
\begin{aligned}
C_1&=p-Q,\\
C_2&=q-Qp+EF+GH,\\
C_3&=EFA-Qq+GHp,\\
C_4&=GHq.
\end{aligned}
\end{equation}
The quartic Routh--Hurwitz conditions are
\begin{equation}
\label{eq:p3-stability}
\begin{gathered}
C_1>0,\quad C_2>0,\quad C_3>0,\quad C_4>0,\\
C_1C_2>C_3,\qquad
C_1C_2C_3>C_3^2+C_1^2C_4.
\end{gathered}
\end{equation}
These inequalities provide a direct analytical or numerical test for local
coexistence stability.

\subsection{Persistence and Bifurcation Behavior}

Local stability describes trajectories close to one equilibrium, whereas
uniform persistence concerns all positive solutions.  The reference combines
the dissipativity result with boundary-flow arguments.  In particular, if
\(P_1\) is globally attracting on the \((S,x)\) boundary subsystem but repels
the missing first predator, and \(P_2\) is globally attracting on the
\((S,x,y)\) boundary subsystem but repels the missing top predator, then no
strictly positive trajectory can approach the boundary.  Under the paper's
additional hyperbolicity and no-cycle assumptions, there is an \(\eta>0\)
such that
\begin{equation}
\liminf_{t\to\infty}S(t),\quad
\liminf_{t\to\infty}x(t),\quad
\liminf_{t\to\infty}y(t),\quad
\liminf_{t\to\infty}z(t)\geq\eta.
\end{equation}
This uniform persistence implies the existence of at least one positive
interior equilibrium \(P_3\).

For bifurcation analysis, the paper introduces a real parameter \(\mu\) into
the prey growth response, writing \(f_1(S;\mu)\), and studies how the Jacobian
eigenvalues change with \(\mu\).  A Hopf bifurcation requires a complex
conjugate pair to cross the imaginary axis with nonzero speed while the other
eigenvalues remain away from it.

At \(P_2\), the cubic factor reaches a candidate Hopf boundary when
\begin{equation}
\label{eq:p2-hopf-boundary}
A_1A_2=A_3,
\end{equation}
with positive coefficients, a nonzero top-predator eigenvalue, and the
transversality condition satisfied.  At \(P_3\), the corresponding quartic
boundary is
\begin{equation}
\label{eq:p3-hopf-boundary}
C_1C_2C_3=C_3^2+C_1^2C_4,
\end{equation}
again together with coefficient positivity and transversality.  Crossing one
of these boundaries can create a family of periodic solutions.  The numerical
oscillatory experiment later in this report illustrates behavior near such a
boundary, but a damped oscillatory trajectory alone is not a proof of Hopf
bifurcation.

\subsection{Consistency of the Implemented Model with the Reference Analysis}

Two inconsistencies in the printed article must be handled explicitly.  First,
the displayed nondimensional equation~(1.2) appears to repeat population
factors inside several per-capita growth terms.  Read literally, that display
does not produce the paper's later Jacobian or equilibrium formulas.  The
system in~\eqref{eq:model} does produce them and is therefore the model used by
the simulator, equilibrium calculations, and this report.  Second, the final
discussion states the opposite sign for washout stability, while Theorem~2.2
and the eigenvalues at \(P_0\) correctly require
\(f_1(1)-D_1<0\).  We follow the theorem and the direct linearization.

The code mirrors the mathematical construction as follows:
\begin{itemize}
\item \texttt{equilibria\_stability/model.py} implements~\eqref{eq:monod},
      its derivative, and its inverse;
\item \texttt{equilibria\_stability/analysis.py} computes \(P_0\) through
      \(P_3\), checks positivity, and evaluates the \(P_0\) threshold;
\item \texttt{demo/app.py} evaluates~\eqref{eq:jacobian} and computes its
      characteristic roots for the interactive local-stability display.
\end{itemize}

Table~\ref{tab:equilibrium-code-check} gives representative outputs produced
by the equilibrium implementation.  They also serve as analytical targets for
the simulations in Section~4.

\begin{table}[htbp]
\centering
\caption{Representative equilibrium calculations from the implemented model.}
\label{tab:equilibrium-code-check}
\small
\begin{tabular}{lll}
\toprule
Scenario & Richest feasible equilibrium & Coordinates \((S,x,y,z)\) \\
\midrule
Washout & \(P_0\) & \((1,0,0,0)\) \\
Prey only & \(P_1\) & \((0.166667,1.666667,0,0)\) \\
Two species & \(P_2\) & \((0.628667,0.333333,0.409333,0)\) \\
Coexistence & \(P_3\) & \((0.055067,0.804597,0.280534,0.081550)\) \\
\bottomrule
\end{tabular}
\end{table}
